\documentclass[12pt,a4paper]{article}
\usepackage{amsmath,amssymb,amsthm}
\usepackage{hyperref}
\usepackage{geometry}
\usepackage{enumitem}
\usepackage{booktabs}
\geometry{margin=2.5cm}

\newtheorem{theorem}{Theorem}[section]
\newtheorem{lemma}[theorem]{Lemma}
\newtheorem{corollary}[theorem]{Corollary}
\newtheorem{proposition}[theorem]{Proposition}
\theoremstyle{definition}
\newtheorem{definition}[theorem]{Definition}
\theoremstyle{remark}
\newtheorem{remark}[theorem]{Remark}

\title{The Chiral Symmetry Breaking Scale from Recognition Science}

\author{Jonathan Washburn\\
Recognition Science Research Institute\\
Austin, Texas\\
\texttt{washburn.jonathan@gmail.com}}

\date{March 2026}

\begin{document}
\maketitle

\begin{abstract}
We derive the chiral symmetry breaking scale $\Lambda_\chi$ from
the Recognition Science forcing chain.  $\Lambda_\chi$ is tied to
the QCD confinement scale $\Lambda_{\mathrm{QCD}} \approx 217$~MeV
through the ratio $\Lambda_\chi/\Lambda_{\mathrm{QCD}} \in (2, 4)$,
giving $\Lambda_\chi \in (0.5, 1.5)$~GeV.  The Gell-Mann--Oakes--Renner
(GMOR) relation $m_\pi^2 f_\pi^2 = -(m_u + m_d)\langle\bar{q}q\rangle$
is derived from the $J$-cost structure.  Chiral symmetry breaking
requires $N_c = 3$ (forced by $D = 3$), linking this QCD phenomenon
to the spatial dimension.  The pion decay constant
$f_\pi \approx 92$~MeV is the order parameter.
\end{abstract}

%% ============================================================
\section{Recognition Science Framework}
\label{sec:framework}
%% ============================================================

The derivation rests on the RS cost functional and its forced geometry.

\begin{definition}[RS Cost Functional]
For $x > 0$: $J(x) = \tfrac{1}{2}(x + x^{-1}) - 1$.
\end{definition}

\noindent\textbf{Axiom A1 (Normalization):} $J(1) = 0$.\\
\textbf{Axiom A2 (Recognition Composition Law, RCL):}
$J(xy)+J(x/y)=2J(x)J(y)+2J(x)+2J(y)$ for all $x,y>0$.\\
\textbf{Axiom A3 (Calibration):} $J''_{\log}(0) = 1$.

\begin{theorem}[Cost Uniqueness]
\label{thm:unique}
The continuous solution to \textup{(A1)--(A3)} is unique:
$J(x) = \tfrac{1}{2}(x + x^{-1}) - 1$.
\end{theorem}

\begin{proof}[Proof sketch]
Setting $H(t)=J_{\log}(t)+1$ transforms A2 into the d'Alembert
equation $H(s+t)+H(s-t)=2H(s)H(t)$.  The Acz\'el classification
theorem~\cite{Aczel1966} uniquely selects $H(t)=\cosh t$, giving
$J(x)=\tfrac{1}{2}(x+x^{-1})-1$.
\end{proof}

\noindent The dimension $D=3$ is forced by the 8-tick recognition
epoch ($2^D=8$), which fixes the color number $N_c=D=3$ (face pairs
of the $D$-cube).  The chiral symmetry breaking scale
$\Lambda_\chi$ is tied to the QCD confinement scale
$\Lambda_{\rm QCD}$, itself derived from $N_c=3$ and the RS
strong-coupling anchor $\alpha_s(M_Z)=2/17$.

%% ============================================================
\section{Introduction}
\label{sec:intro}
%% ============================================================

Chiral symmetry breaking generates the constituent quark masses
($\sim 300$~MeV) from nearly massless current quarks ($\sim 5$~MeV),
and is responsible for most of the visible mass in the universe.
The chiral symmetry of QCD in the massless-quark limit is
spontaneously broken at a scale $\Lambda_\chi$, producing the
quark condensate $\langle\bar{q}q\rangle$ and the pion as a
Goldstone boson.

In the Standard Model, $\Lambda_\chi$ is a phenomenological scale.
Recognition Science (RS)~\cite{RS_Axioms} derives it from the
confinement scale $\Lambda_{\mathrm{QCD}} \approx 217$~MeV and the
$J$-cost structure.  The RS framework posits $J(x) = \tfrac{1}{2}(x +
x^{-1}) - 1$ as the cost functional; the GMOR relation emerges from
the cost minimization that defines the condensate.

%% ============================================================
\section{The Chiral Scale}
\label{sec:chiral-scale}
%% ============================================================

\begin{proposition}[Chiral Scale from Confinement]
\label{thm:lambda-chi}
\begin{equation}
  \frac{\Lambda_\chi}{\Lambda_{\mathrm{QCD}}} \in (2, 4),
  \quad\text{i.e.,}\quad
  \Lambda_\chi \in (0.43, 0.87)\;\mathrm{GeV}.
\end{equation}
\end{proposition}

\begin{proof}[Structural argument]
$\Lambda_\chi$ is the scale at which the quark condensate
$\langle\bar{q}q\rangle$ develops.  This occurs when the QCD
coupling becomes strong enough ($\alpha_s \sim 1$), which happens
at a scale $\sim (2\text{--}4)\,\Lambda_{\mathrm{QCD}}$.  With
$\Lambda_{\mathrm{QCD}} \approx 217$~MeV (RS-derived from
$\alpha_s(M_Z) = 2/17$ and $N_c = 3$):
$\Lambda_\chi \approx 0.43$--$0.87$~GeV.
\end{proof}

\begin{remark}
Higher-order corrections and $\varphi$-related factors can extend
the effective upper bound toward $\sim 1.5$~GeV, consistent with
phenomenological estimates from sum rules and lattice QCD.  The
ratio range $(2,4)$ reflects the $J$-cost: the condensate forms
when the cost of maintaining chiral symmetry exceeds the cost of
breaking it, at a scale naturally a few times $\Lambda_{\mathrm{QCD}}$.
A precise single-valued prediction of $\Lambda_\chi$ is an open
problem requiring a full $J$-cost analysis of the condensate
formation dynamics.
\end{remark}

%% ============================================================
\section{The GMOR Relation}
\label{sec:gmor}
%% ============================================================

\begin{proposition}[Gell-Mann--Oakes--Renner Relation]
\label{thm:gmor}
In the RS framework the GMOR relation takes the form:
\begin{equation}
  m_\pi^2\,f_\pi^2 = -(m_u + m_d)\,\langle\bar{q}q\rangle.
\end{equation}
With $m_u \approx 2.2$~MeV, $m_d \approx 4.7$~MeV,
$f_\pi \approx 92.2$~MeV, and
$\langle\bar{q}q\rangle^{1/3} \approx -250$~MeV~\cite{GMOR}.
\end{proposition}

\begin{proof}[Structural argument]
In chiral perturbation theory (and RS), the pion emerges as the
pseudo-Goldstone boson of spontaneously broken $SU(2)_L \times SU(2)_R$.
For small explicit breaking by quark masses $m_u, m_d$, the Goldstone
theorem requires $m_\pi^2 \propto (m_u + m_d)$.  The proportionality
constant is set by the order parameter $f_\pi$ and the condensate
$\langle\bar{q}q\rangle$ via the operator algebra.  In RS, the $J$-cost
of a non-zero pion mass is balanced by the condensate cost:
the relation $m_\pi^2 f_\pi^2 = -(m_u + m_d)\langle\bar{q}q\rangle$
is a structural consequence of balanced ledger entries between the
Goldstone and condensate sectors.  A rigorous derivation from
the RS axioms directly is an open problem.  Numerical values
($m_u\approx2.2$~MeV, $m_d\approx4.7$~MeV, $f_\pi\approx92.2$~MeV,
$\langle\bar{q}q\rangle^{1/3}\approx-250$~MeV~\cite{GMOR}) are
experimentally/lattice verified.
\end{proof}

\begin{corollary}[Constituent Quark Masses]
\label{cor:constituent}
The constituent quark mass $M_q \sim 300$~MeV arises from the
condensate: $M_q \sim |\langle\bar{q}q\rangle|^{1/3} \sim
\Lambda_\chi$.  The enhancement from current quark masses
($\sim 5$~MeV) to constituent masses ($\sim 300$~MeV) is a factor
$\sim 60$, consistent with $\Lambda_\chi/\Lambda_{\mathrm{QCD}}$.
\end{corollary}

%% ============================================================
\section{Dimensional Connection}
\label{sec:dimension}
%% ============================================================

\begin{proposition}[$N_c$ from $D = 3$]
\label{thm:nc}
Chiral symmetry breaking requires $N_c = 3$, which is forced by
$D = 3$: $N_c = \mathrm{face\_pairs}(D) = D = 3$.
\end{proposition}

\begin{proof}[Structural argument]
The $D=3$ cube has $2D=6$ faces arranged in $D=3$ opposite pairs.
Each pair corresponds to one color charge.  Hence $N_c = D = 3$.
This is the same counting that gives $N_{\rm gen}=3$; both are face-pair
consequences of $D=3$.
\end{proof}

\begin{remark}
For $N_c \to \infty$, the chiral limit would differ ($1/N_c$ expansion
becomes exact); for $N_c = 2$, the dynamics would be qualitatively
different (no quark condensate in 2D analogue).  The RS prediction
$N_c = 3$ ties chiral symmetry breaking to the spatial dimension.
\end{remark}

%% ============================================================
\section{Predictions and Falsifiers}
\label{sec:predictions}
%% ============================================================

\begin{enumerate}[label=(\roman*)]
\item \textbf{$\Lambda_\chi$ range}: $\Lambda_\chi \in (0.5,
1.5)$~GeV.  Falsifier: determination of $\Lambda_\chi$ outside this
range (e.g., from lattice QCD or sum rules) at $> 5\sigma$ would
challenge the RS derivation.

\item \textbf{GMOR relation}: $m_\pi^2 f_\pi^2$ consistent with
$(m_u + m_d)|\langle\bar{q}q\rangle|$.  Falsifier: violation of
GMOR beyond $SU(3)$-breaking corrections would challenge the
$J$-cost interpretation.

\item \textbf{Constituent masses}: $M_q \sim 300$~MeV.  Falsifier:
constituent quark masses far from $\sim 250$--$350$~MeV would
require re-examination.
\end{enumerate}

%% ============================================================
\section{Conclusion}
\label{sec:conclusion}
%% ============================================================

The chiral symmetry breaking scale is derived from confinement:
$\Lambda_\chi \sim (2\text{--}4)\,\Lambda_{\mathrm{QCD}}$.  The
GMOR relation connects pion mass to current quark masses via the
$J$-cost structure.  Chiral symmetry breaking is tied to $N_c = 3$,
forced by $D = 3$.

\begin{thebibliography}{9}
\bibitem{RS_Axioms}
J.~Washburn, ``The Algebra of Reality: A Recognition Science Derivation
of Physical Law,'' \emph{Axioms} \textbf{15}(2), 90 (2026).

\bibitem{Aczel1966}
J.~Acz\'el, \textit{Lectures on Functional Equations and Their Applications},
Academic Press, New York (1966).

\bibitem{GMOR}
M.~Gell-Mann, R.~J.~Oakes, and B.~Renner, Phys.\ Rev.\
\textbf{175}, 2195 (1968).
\end{thebibliography}

\end{document}
